\chapter{Introduction}
\label{ch:introduction}

\section{Background and Motivation}
\label{sec:1_1}

The proliferation of online marketplaces has fundamentally transformed how consumers search for and secure housing. In Switzerland alone, platforms such as Homegate and ImmoScout24, operated by Swiss Marketplace Group (SMG), facilitate hundreds of thousands of real estate transactions annually. However, this digital transformation has simultaneously created fertile ground for sophisticated fraud schemes, generating complex data environments where fraudulent patterns are embedded within massive transaction graphs.

Real estate fraud in online marketplaces manifests in several forms: phantom listings advertising non-existent properties, identity theft schemes harvesting applicant information, and coordinated fraud rings operating networks of fake accounts. These fraud operations leave distinctive traces in marketplace data shared device fingerprints, IP address clusters, temporal posting patterns, and entity relationship networks that are invisible when analyzing listings in isolation but become apparent through systematic data mining.

Traditional fraud detection approaches treat listings as independent observations, applying rule-based systems or commercial services such as SEON that analyze individual transactions without exploiting relational structure. While these approaches capture obvious fraud indicators (disposable emails, high-risk IP countries), they cannot discover the network patterns that distinguish fraud rings from legitimate property management operations.

Despite the maturity of graph-based fraud detection in adjacent domains---with over 100 publications applying Graph Neural Networks to financial transactions \cite{cheng2024gnn}, production deployments at e-commerce platforms processing billions of transactions (eFraudCom at Taobao, C-FATH at JD.com), and standardized benchmark datasets (YelpChi, Elliptic)---no academic research has applied these methods to real estate marketplace fraud. A systematic review of fraud detection literature reveals this absence: Mutemi and Bação's \cite{mutemi2024ecommerce} survey of 101 e-commerce fraud publications excludes real estate entirely; Cherif et al.'s \cite{cherif2024financial} systematic review of GNN-based financial fraud detection covers credit cards, money laundering, and insurance but omits property marketplaces; the only peer-reviewed real estate fraud study \cite{mohdamin2024clustering} employs traditional clustering rather than graph methods. This gap persists despite documented fraud losses exceeding \$145 million annually in real estate scams \cite{fbi2023ic3} and the inherent relational structure of marketplace data---agent networks, cross-platform listing patterns, device-sharing graphs---that suggests graph methods would be highly applicable.

This thesis addresses this gap by applying graph-based data mining techniques to discover, characterize, and validate fraud patterns in real estate marketplace data. By constructing heterogeneous graphs from entity relationships and applying graph-based pattern extraction, we aim to uncover relational fraud signatures invisible to attribute-only analysis, while investigating the unique structural properties of fraud in this unexplored domain.

\section{Research Objectives and Questions}
\label{sec:1_2}

The primary objective of this thesis is to discover, characterize, and validate relational fraud patterns in real estate marketplace data through graph-based data mining, demonstrating their utility for improved fraud detection and operational investigation support.

This objective is operationalized through four research questions:

\textbf{RQ1: What relational fraud patterns exist in marketplace graphs, and do they provide predictive signal beyond individual listing attributes?}

This question addresses pattern discovery and validation. By constructing heterogeneous graphs from entity relationships and extracting relational representations, we identify structural patterns distinguishing fraudulent from legitimate listings. Comparison between graph-augmented and attribute-only models validates whether discovered patterns carry predictive information.

\textbf{RQ2: What entity-sharing patterns characterize fraud rings versus legitimate multi-listing users?}

Beyond aggregate pattern detection, this question examines the specific topological and behavioral signatures of coordinated fraud operations. We analyze device-sharing networks, IP clustering, email/phone reuse patterns, and temporal coordination to characterize fraud ring structure versus legitimate property management behavior.

\textbf{RQ3: How do fraud patterns evolve temporally, and what data requirements enable effective pattern discovery?}

Fraud patterns are non-stationary; fraudsters adapt tactics over time. This question examines temporal evolution of fraud signatures and, critically, identifies the minimum data scale at which graph-based patterns emerge as predictive. Through expanding window analysis, we quantify data requirements for effective pattern mining.

\textbf{RQ4: What actionable fraud indicators emerge from pattern analysis that can support human investigation?}

Discovered patterns must translate to operational value. This question examines which features both individual attributes and relational patterns most strongly indicate fraud, producing interpretable indicators that fraud analysts can use to prioritize and guide investigations.

\section{Scope and Delimitations}
\label{sec:1_3}

This research is conducted as a case study utilizing production data from Swiss Marketplace Group's real estate platforms.

\textbf{In Scope:}
\begin{itemize}
    \item Exploratory analysis of marketplace data structure and fraud distributions
    \item Construction of heterogeneous entity-relationship graphs from listing data
    \item Graph-based pattern extraction using Graph Neural Network techniques
    \item Validation of pattern utility through predictive modeling comparison
    \item Temporal pattern analysis through expanding window experiments
    \item Interpretability analysis to identify actionable fraud indicators
\end{itemize}

\textbf{Out of Scope:}
\begin{itemize}
    \item Real-time streaming analysis (batch processing focus)
    \item Unstructured content mining (listing text, images)
    \item Cross-marketplace pattern transfer
    \item Fully dynamic graph representations
\end{itemize}

\textbf{Delimitations:}
\begin{itemize}
    \item Patterns reflect Swiss real estate market characteristics; generalization requires validation
    \item Ground truth derives from operational flagging, which may contain labeling noise
    \item Proprietary data constraints limit reproducibility; statistical summaries provided
\end{itemize}

\section{Research Contributions}
\label{sec:1_4}

This thesis makes contributions to both the data mining research literature and operational fraud detection practice at Swiss Marketplace Group.

\subsection*{Academic Contributions}

\textbf{Contribution 1: Discovery of the Heterophily Paradox in Real Estate Fraud Graphs.}
We discover that real estate marketplace graphs exhibit \emph{inverted} fraud topology compared to assumptions in the graph fraud detection literature. While existing work assumes fraudsters form densely connected clusters requiring camouflage-resistant methods \cite{dou2020enhancing,liu2020alleviating}, our analysis reveals the opposite pattern: legitimate property managers form dense entity-sharing networks (5.3\% fraud rate for devices with 5+ listing connections), while fraudsters operate in relative isolation (16.7\% fraud rate for single-listing devices). This heterophily pattern---where high connectivity signals legitimacy rather than fraud---explains why standard GNN architectures underperform in this domain and represents the first empirical characterization of fraud network structure in real estate marketplaces.

\textbf{Contribution 2: Establishment of Graph Density Threshold for Relational Pattern Mining.}
Through systematic expanding window analysis, we provide the first empirical quantification of data requirements for graph-based fraud pattern discovery. Graph mining provides predictive lift over attribute-only analysis only when training data exceeds approximately 300,000 samples---the threshold at which fraud ring connectivity reaches sufficient density to overcome legitimate power-user noise. Below this threshold, relational patterns lack structural coherence for meaningful extraction. This finding addresses a methodological gap in the literature: prior work demonstrates graph patterns improve detection but does not investigate \emph{when} this occurs, leaving practitioners unable to assess graph-based approach viability for their data contexts.

\textbf{Contribution 3: Demonstration of Temporal Robustness Inversion Under Production Evaluation.}
We demonstrate that evaluation methodology fundamentally determines model selection conclusions in fraud detection. Under fixed-split evaluation (standard academic benchmark), graph-augmented models achieve +5.3\% AUC-PR improvement over tabular baselines. However, under expanding window evaluation (production deployment scenario with periodic retraining), this relationship inverts: tabular models outperform graph-augmented approaches by 3$\times$ (0.63 vs. 0.22 AUC-PR). This discovery reveals that fraud evolution operates at the \emph{topology level}---new fraud campaigns form disconnected graph components with distinct structures---not merely at the feature distribution level assumed in prior concept drift literature \cite{gama2014survey}. We identify two operational regimes: forensic analysis (offline, large historical graphs) where GNN methods excel, and real-time detection (incremental retraining) where tabular methods provide superior robustness.

\subsection*{Practical Contributions}

\textbf{Contribution 4: Complementary Detection Layer for Existing Fraud Systems.}
The discovered patterns enable a secondary detection mechanism operating alongside SMG's existing SEON-based fraud detection. Unlike SEON's rule-based approach analyzing individual transactions, the graph-based system identifies fraud through relational signatures---detecting coordinated activity invisible to single-listing analysis. This provides both forward detection (risk scoring at submission) and backward detection (retrospective identification of published fraud based on emerging entity-sharing patterns from subsequent user activity).

\textbf{Contribution 5: Empirically-Validated Rule Enhancement.}
The characterized patterns---super-connector thresholds (10+ listings per device yields 35.5\% fraud rate, 4.4$\times$ baseline), cross-platform signatures (dual Homegate/ImmoScout24 posting yields 26.74\% fraud rate), temporal indicators (6-7 AM submission peaks), and geographic concentration (West African origin patterns showing 60-90\% fraud rates)---provide empirically quantified signals for SEON rule enhancement. Each pattern includes precision metrics enabling risk-calibrated rule implementation rather than heuristic thresholds.

\section{Thesis Organization}
\label{sec:1_5}

The remainder of this thesis is organized as follows:

\textbf{Chapter 2: Literature Review} surveys related work in graph-based pattern mining, anomaly detection in networks, fraud detection methodologies, and interpretable machine learning, positioning this research within the data mining literature.

\textbf{Chapter 3: Methodology} details the data mining approach including exploratory data analysis, graph construction methodology, pattern extraction techniques, validation experimental design, and temporal analysis protocols.

\textbf{Chapter 4: Implementation} describes technical realization including data pipelines, graph construction, pattern extraction via Graph Neural Networks, and pattern validation.

\textbf{Chapter 5: Results and Analysis} presents discovered patterns organized by research question: relational structures (RQ1), fraud ring characteristics (RQ2), temporal evolution and data requirements (RQ3), and actionable indicators (RQ4). The chapter concludes with interpretation of findings and practical recommendations.

\textbf{Chapter 6: Conclusion and Future Work} summarizes contributions, examines limitations, and identifies directions for extended pattern discovery research.
